\section{Polynomial convolution over arbitrary fields}\label{sec:pack}

Multiplying two polynomials of degree less than $n$ amounts to
computing the convolution of their two coefficient lists. If the field
contains suitable roots of unity, one can evaluate both polynomials,
multiply their values pointwise and interpolate. We will obtain the same
logarithmic exponent over fields that may lack those roots.

\begin{theorem}[Formal polynomial multiplication]\label{thm:conv}
For every $\alpha>1/2$, a division-free circuit over $K$ computes the
$2n-1$ coefficients of the product of two length-$n$ polynomials in
$O_\alpha(n(1+\log n)^\alpha)$ binary additions, subtractions and
multiplications. The count includes multiplication by prescribed scalars
and the implied constant is uniform over all fields $K$.
\end{theorem}

Each scalar in a root extension of degree $e$ occupies $e$ base-field
coordinates. We use each coordinate to hold input coefficients or
their products. The computation is a convolution on $N$ cyclic
addresses, carrying $r$ extension scalars per address as the
coefficients of a truncated polynomial.
Its total dimension $erN$ is comparable to the original length $n$.
Multiplying at one address reduces to ordinary convolution over $K$
at a length $O(er)$ whose logarithm is much smaller than $\log n$.

The choice of $N$ supplies a grid of small prime axes. In positive
characteristic the extension contains the required roots on every axis.
In characteristic zero we adjoin roots of order $p-1$ on an axis of
prime length $p$. A change of basis and a polynomial deformation then
give a pointwise multiplication using these shorter transforms.
Cantor and Kaltofen gave a convolution algorithm over arbitrary
algebras using $O(n\log n\log\log n)$ operations
\cite{CantorKaltofen}. For smooth extension degrees and many-axis
transforms in polynomial arithmetic see
\cite{harveyhoevenfinite,harveyhoevenconditional}.

\subsection{Packing the product without wraparound}

Suppose $E=K(\xi)$ has degree $e$, and choose positive integers $r,N$
with $8n\le erN\le9n$. Write $J_h=E[z]/(z^h)$ for truncated
polynomials of length $h$ (jets). In $J_r[X]/(X^N-1)$ the monomials
$\xi^a z^b X^j$, $0\le a<e$, $0\le b<r$, $0\le j<N$, form a
$K$-basis. To pack the inputs, reserve half of the room on each axis
for the growth in exponents when two monomials are multiplied.
Take radices
\[
 E_0=\lceil e/2\rceil,\qquad R=\lceil r/2\rceil,\qquad
 U=\lceil N/2\rceil.
\]
Since $E_0 RU\ge erN/8\ge n$, each input fits in the lower part of this
array. With $0\le a<E_0$, $0\le b<R$ and $0\le j<U$, place its coefficients
by the rule
\begin{equation}\label{poly:eq:mixed-radix-packing}
           a+E_0 b+E_0 Rj \quad\longmapsto\quad \xi^a z^b X^j,
\end{equation}
padding unused positions with zeros. The inequality
$2(\lceil k/2\rceil-1)<k$ applied at $k=e,r,N$ shows that
multiplying two occupied monomials adds their exponents without
reduction by the minimal polynomial of $\xi$, by $z^r$, or by $X^N-1$.

In the product, gather every coordinate $(a,b,j)$ with
$0\le a<e$, $0\le b<r$ and $0\le j<N$ by the integer $a+E_0 b+E_0 Rj$, retaining
indices $0,\ldots,2n-2$. Different triples can contribute to the same
index because of carries in the smaller radices. The index of each
product monomial is the sum of its two input indices, so gathering
gives ordinary convolution at all $2n-1$ outputs. Packing and
gathering cost $O(erN)$.

\subsection{Choosing the cyclic length and the field extension}

Put $T=\lceil\log n\rceil$. There is an absolute $C$ such that, with
$B=CT$ and for all sufficiently large $T$, there is a list of distinct odd primes
$p_i\le B$ whose product exceeds $n^2$ and
\[
 \log\mathrm{lcm}_i(p_i-1)=O(B^{4/5}\log B).
\]
We may omit any one forbidden prime. Shifted smooth primes have a similar role in matrix nonrigidity
\cite[Section~4]{DvirLiu}.
To see the claim, count primes $p\in[B/2,B]$
with $p\equiv1\bmod q$, summing over primes $B^{1/5}\le q\le B^{1/4}$.
By the Bombieri--Vinogradov theorem
\cite{BombieriV}, the sum is
\[
  (\operatorname{Li}(B)-\operatorname{Li}(B/2))
              \sum_q 1/(q-1)+O(B/(\log B)^2)
   = \Omega(B/\log B)
\]
by Mertens' theorem. Each $p$ is counted at most five times.
All prime factors of each $p-1$ are at most $B^{4/5}$, so
there are at most $B^{4/5}$ prime powers (largest power for each prime)
to include in the least common multiple, and each is at most $B$.
Taking $C$ large gives the claim.

Let $\ell=2$ except in characteristic two, where we take $\ell=3$.
This makes the radix invertible in $K$. In positive characteristic
also omit the characteristic from the list of primes.
In characteristic zero take primitive roots of all orders $p_i-1$,
and in positive characteristic take primitive roots of all orders $p_i$.
In both cases include a primitive root at the largest $\ell$-power
at most $8B$. Its powers supply the smaller radix lengths.
Let $E=K(\xi)$, where $\xi$ is a primitive root of order the least
common multiple of all these orders, and put $e=[E:K]$.

We claim that $\log e=O(B^{4/5}\log B)$. In characteristic zero
the order of $\xi$ bounds $e$, so the claim follows from the estimate
on the least common multiple of the $p_i-1$. In characteristic $b>0$,
we can choose $\xi\in\mathbb F_{b^D}$ for
\[
 D= \mathrm{lcm}_i(p_i-1)\cdot\varphi(\text{largest $\ell$-power up to $8B$}).
\]
Indeed $b$ is a unit modulo every required order. Its order modulo
$p_i$ divides $p_i-1$, and its order modulo any required $\ell$-power
divides the indicated totient. Hence $b^D-1$ is divisible by all the
orders. Since $E\subseteq K\mathbb F_{b^D}$, we have $e\le D$, as
required.

Take $N=\prod_{i=1}^d p_i$ from the longest prefix with
$eN\le n/B$.
Since $e=2^{o(T)}$ while the full product exceeds $n^2$, the
prefix is nonempty and has a next prime for all large $T$. That prime
is at most $B$ and would violate the upper bound, so
$n/B^2<eN\le n/B$. Choose an integer $r$ so that
\[
                8n\le erN\le 9n.
\]
The number $d$ of axes is at most $B$, and $B\le r\le9B^2$. In
particular
\[
             erN=\Theta(n),\qquad \log(6er)=O(T^{4/5}\log T).
\]
The first estimate accounts for the storage used by the packing above.
The second provides the saving in the recursive calls.

\subsection{Counting gates over the base field}

Let $M_K(j)$ denote the minimum circuit size of ordinary length-$j$
polynomial convolution over $K$. To charge a family of shorter
convolutions by its total operand length, put
$\mu_K(s)=1+\max_{1\le j\le s}M_K(j)/j$.

\begin{lemma}[Multiplication and shifts in jet coordinates]\label{poly:lem:descent-jet-cost}
Multiplication in $J_h$ costs $O(eh\mu_{K}(2eh))$ gates.
In characteristic zero, for any fixed $c\in K$, the shift $P(z)\mapsto P(z+c)$ of a
degree-$<h$ polynomial over $E$ costs $O(eh\mu_{K}(h))$ gates.
\end{lemma}
\begin{proof}
Represent the two jet inputs by bivariate polynomials of degrees $<e$
in $Y$ and $<h$ in $z$. Substitute $Y=w$ and $z=w^{2e-1}$. Raw product
exponents in $Y$ are at most $2e-2$, so the substitution separates the
rows of the product. Each packed operand has length at most $2eh$.
Multiply as ordinary polynomials over $K$. Keep the first $h$ rows and reduce each modulo the fixed monic minimal
polynomial $D$ of $\xi$ over $K$. For a row $V=DQ+S$ of degree $<2e$
with $\deg S<e$, reversing at degrees $2e-1,e,e-1$ gives
\begin{equation}\label{poly:eq:reciprocal-reduction}
  V^{\rm rev}\equiv D^{\rm rev}Q^{\rm rev}\pmod {Y^e}.
\end{equation}
Truncate $V^{\rm rev}$ modulo $Y^e$, multiply there by the precomputed
inverse of $D^{\rm rev}$, reverse the quotient, and subtract $DQ$. The two products per row
have operand lengths at most $e+1\le2e$. This reduction uses only
monicity and works over any coefficient ring. Together with the packed
product, all these multiplications have total operand length $O(eh)$.

For the shift, the new coefficient of $z^b$ is
\[
  \frac1{b!}\sum_{j=b}^{h-1} j!\,[z^j]P\,
                             \frac{c^{j-b}}{(j-b)!}.
\]
After reversal this is one fixed convolution of length $h$ per
$K$-coordinate, together with fixed scalings.
\end{proof}

A convolution with a fixed operand is a linear map at the same
asymptotic cost, even when the available multiplication circuit is an
arbitrary division-free circuit. Retain the constant and linear parts
of each wire in the remaining input. At a product gate the linear part
is $a_0b_1+a_1b_0$, where $a_0,b_0$ are precomputed constants.
The linear output is the required map.

Consider one layer of Section~\ref{sec:lay} on $V=O(N)$ registers
with entries in $J_h$. A shear block of size $2^j$ acts independently
on the $eh$ base-field coordinates of each entry. Since $j=O(1+T)$,
Lemma~\ref{lem:shear} bounds its cost by
$O_\beta(2^j eh(1+T)^\beta)$ for $\beta>1/2$. A monomial layer needs
at most $V$ multiplications by prescribed jet scalars. Summing the block sizes for the shear layer and applying
Lemma~\ref{poly:lem:descent-jet-cost} to the monomial layer bounds
the cost in either case by
\begin{equation}\label{poly:eq:descent-sync}
 O_\beta\big(Veh((1+T)^\beta+\mu_K(2eh))\big).
\end{equation}
For a tensor product of Fourier matrices of sides $a\le B$, suppose
their orders are invertible and their roots are in $E$. On each axis
take the least $\ell$-power at least $2a$, which is less than $6a$.
The roots for this auxiliary length are in $E$ by construction.
Radix stages and Lemma~\ref{up:toeplitz} give a word of
$O(1+\log B)$ layers on each axis with a bounded ratio of workspace
to data. Lemma~\ref{up:fibers} synchronizes the axis words in the same
depth and at constant total workspace overhead. Thus the cost of the
whole transform is Equation~\eqref{poly:eq:descent-sync} times
$O(1+\log B)$.

\paragraph{Positive characteristic.}
Chinese remaindering identifies the cyclic $N$-address space with a
$p_1\times\cdots\times p_d$ grid under componentwise addition.
Every axis order is invertible and $E$ contains a primitive root of
each order. Two forward transforms, $N$ products in $J_r$, and one
inverse transform therefore compute cyclic convolution. Fix
$1/2<\beta<\alpha$. Since $erN=\Theta(n)$, this costs at most
\begin{equation}\label{poly:eq:descent-cost}
 O_\beta\big(n(1+\log B)((1+T)^\beta+\mu_K(6er))\big).
\end{equation}
The argument in characteristic zero will allow jets of length $3r$,
which accounts for the factor $6$ here.

\subsection{A binomial coordinate on a prime axis in characteristic zero}

We now pass from the prime axis of length $p$ to an axis of length
$p-1$, whose roots lie in $E$.
Fix a prime $p$ from the list, put $m=p-1$, and choose a primitive
$m$th root $\omega\in E$. Let
$A=E[x]/(1+x+\cdots+x^{p-1})$.

\begin{lemma}[Prime cyclic algebra in binomial coordinates]\label{poly:lem:descent-prime}
There is a $\lambda\in E^\times$ and an isomorphism
\[
 E[x]/(x^p-1)\simeq E\times E[Y]/(Y^m-\lambda).
\]
The coordinate change in either direction has a word of
$O(1+\log p)$ layers on its $p$ coordinates with $O(p)$ restored workspace.
\end{lemma}
\begin{proof}
The factors $x-1$ and $1+x+\cdots+x^{p-1}$ are coprime since the
second takes the nonzero value $p$ at $1$. Chinese remaindering
gives a scalar coordinate by evaluation at $1$ and a second coordinate
in $A$. Choose a generator $g$ of $(\Z/p\Z)^\times$.
Order the nonzero residues by powers of $g$, as in Rader's
prime-length Fourier transform~\cite{rader}. The powers $x^{g^t}$ for
$0\le t<m$, with exponents taken modulo $p$, form a basis of $A$. They span since their sum is $-1$, and there
are $m=\dim_E A$ of them. The automorphism $x\mapsto x^g$ permutes
them cyclically. Its Fourier eigenbasis consists of the Lagrange resolvents
\[
 \gamma_j=\sum_{t=0}^{m-1}\omega^{-jt}x^{g^t},
 \qquad 0\le j<m,
\]
with eigenvalues $\omega^j$ and $\gamma_0=-1$. The algebra $A$ is
reduced because its defining polynomial has simple roots. Hence the
nonzero element $\gamma=\gamma_1$ has nonzero powers. In particular
$\gamma^m$ belongs to the invariant line $E\cdot1$ and equals a
nonzero scalar $\lambda\in E^\times$. Thus $\gamma$ is a unit even
when $A$ is a product of fields. For $1\le j<m$ we have
$\gamma^j=\kappa_j\gamma_j$ for some $\kappa_j\in E^\times$.
Thus $Y\mapsto\gamma$ identifies the binomial quotient with $A$.

In the cyclic algebra write an input as
$v_*+\sum_t v_t x^{g^t}$ and set $w_j=\sum_t\omega^{jt}v_t$.
Fourier inversion gives its image in $A$ as
$v_*+m^{-1}\sum_j w_j\gamma_j$. The coordinate in the scalar factor,
the constant coefficient in the binomial factor, and its positive-degree
coefficients are respectively
\[
  v_*+w_0,\qquad v_*-w_0/m,\qquad
  \frac{w_j}{m\kappa_j}\quad(1\le j<m).
\]
Lemma~\ref{up:toeplitz}, using the least $\ell$-power at least $2m$ as working length, applies the transform on its $m$
coordinates with restored workspace while leaving $v_*$ in place. The remaining two-coordinate
map has determinant $-p/m$, so it and its inverse have constant-length
words by elementary elimination and the shear construction.
The other steps are permutations and nonzero scalings.
The inverse transform is a forward transform at $\omega^{-1}$ followed
by scaling by $1/m$.
\end{proof}

For $p=3$ the nonconstant eigenvector is $\gamma=x-x^2$. The relation
$x^2+x+1=0$ gives $\gamma^2=-3$, so the binomial factor in this case
is $E[Y]/(Y^2+3)$.

The Chinese remainder isomorphism on the address group sends $X^j$
to the monomial of exponents $(j\bmod p_\nu)_\nu$. It identifies
$J_r[X]/(X^N-1)$ with the tensor product of the prime cyclic algebras
over $J_r$ by a permutation of coefficients. Apply
Lemma~\ref{poly:lem:descent-prime} on all axes and distribute this tensor
product over the scalar and binomial factors. The resulting blocks
are indexed by subsets $S\subseteq\{1,\ldots,d\}$. Writing
$m_\nu=p_\nu-1$, the block for $S$ is
\[
 \mathcal T_S=J_r[Y_\nu:\nu\in S]/
             (Y_\nu^{m_\nu}-\lambda_\nu:\nu\in S).
\]
Its width is $W_S=\prod_{\nu\in S}m_\nu$, and
$\sum_S W_S=\prod_\nu(1+m_\nu)=N$. By
\eqref{poly:eq:descent-sync} and Lemma~\ref{up:fibers}, both input conversions and the final inverse
conversion fit within \eqref{poly:eq:descent-cost}.

\subsection{Multiplying a block with one deformation parameter}
Fourier evaluation would multiply each block pointwise if its defining
binomials split, but $\lambda_\nu$ need not have an $m_\nu$th root in $E$.
We connect the block to one with relations $Y_\nu^{m_\nu}=1$ by varying a
single parameter $t$. We compute the product in a truncation at $t=0$
long enough to recover its coefficients as polynomials in $t$, then take
their Taylor expansions at $t=1$. Li, Moreno Maza and Schost used a common parameter
to deform relations in separate variables~\cite[Section~2.2, proof of Theorem~2]{triangularLMS}.
Bostan, Chowdhury, van der Hoeven and Schost developed the method for
triangular sets~\cite[Section~4.2, Proposition~7]{homotopyBCHS}.
Here the inputs themselves carry $r$ coefficients in the parameter,
and each binomial relation adds at most one to the degree required
for recovery.

\begin{lemma}[Multiplication through a split block]\label{poly:lem:descent-deform}
Let $H=2r+d$. Multiplication in $\mathcal T_S$ reduces to a tensor
Fourier transform on each of two arrays over $J_H$, one inverse tensor
transform, $W_S$ pointwise products in $J_H$, and a constant number
of jet scalings and polynomial shifts per basis monomial. The Fourier
axes have widths $m_\nu$, $\nu\in S$.
\end{lemma}
\begin{proof}
Use the $E[t]$-algebra
\[
 \mathcal A=E[t][Y_\nu:\nu\in S]/
 \bigl(Y_\nu^{m_\nu}-f_\nu(t):\nu\in S\bigr),
 \qquad f_\nu(t)=1+(\lambda_\nu-1)t.
\]
Its basis over $E[t]$ consists of the monomials whose exponent on each axis
is less than $m_\nu$, by successive division by the separate monic
relations. Use this same basis after each specialization. For the
$p=3$ example above the relation is $Y^2=1-4t$. In the product
$(a+bY)(c+dY)$ of two constant lifts, the scalar coefficient is
$ac+bd(1-4t)$. A single wrap changes its value from $ac+bd$ at $0$
to $ac-3bd$ at $1$.

Use $z_1=t-1$ for the parameter in $J_r$ at $t=1$, and $z_0=t$
for the parameter in $J_H$ at $t=0$. At the first endpoint, the map
from the specialized algebra to $\mathcal T_S$ sends the class of
$Y_\nu$ to $u_\nu Y_\nu$, where
$u_\nu^{m_\nu}=f_\nu(1+z_1)/\lambda_\nu$.
At the second endpoint the target is the split block with
$Y_\nu^{m_\nu}=1$. The corresponding map sends the class of $Y_\nu$
to $v_\nu Y_\nu$, where $v_\nu^{m_\nu}=f_\nu(z_0)$.
Each right-hand side has constant term one. The root with constant term
one exists and is unique, since the next unknown coefficient occurs
with multiplier $m_\nu\in E^\times$. These assignments respect the
defining relations and multiply each basis monomial by a unit.
We precompute that multiplier and its inverse for each monomial.

Start with two elements of $\mathcal T_S$ and undo the isomorphism
at $t=1$. For each coefficient choose its degree-$<r$ representative
in $z_1$ and substitute $t-1$ for $z_1$. These coefficients specify lifts
of the two elements to $\mathcal A$. In a product of basis monomials
$Y^aY^b$, the coefficient introduced by reduction is
\[
  \prod_{\nu\in S:\ a_\nu+b_\nu\ge m_\nu} f_\nu(t).
\]
Each axis contributes at most one factor, so the coefficient polynomials
of the lifted product have degree at most $2r-2+|S|<H$. Consequently their residues modulo $t^H$ give every coefficient of
the actual product polynomial. Reading these representatives in $E[t]$
lets us evaluate at $t=1+z_1$ even though the local calculation used
only jet arithmetic at zero (Figure~\ref{poly:fig:deformation}).
\begin{figure}[!htbp]
\centering
\begin{tikzpicture}[
  box/.style={font=\small,align=center,text width=4.2cm,minimum height=1.04cm,
              draw=#1,rounded corners=2pt,inner sep=3pt,fill=#1!4},
  labeltext/.style={font=\footnotesize,align=center,inner sep=3pt},
  arrow/.style={-{Stealth[length=1.8mm]},draw=ink,line width=.55pt}]
\node[box=sea] (in) at (0,0)
  {Inputs in $\mathcal T_S$\\$Y_\nu^{m_\nu}=\lambda_\nu$};
\node[box=sea] (out) at (8.4,0)
  {Output in $\mathcal T_S$\\$Y_\nu^{m_\nu}=\lambda_\nu$};
\node[box=ink] (lift) at (0,-2.23)
  {Lifted coefficients in $E[t]$\\$\deg P_a,\deg Q_a<r$};
\node[box=ink] (recover) at (8.4,-2.23)
  {Product coefficients in $E[t]$\\$\deg R_a\le2r-2+|S|$};
\node[box=rust] (splitin) at (0,-4.46)
  {Inputs over $E[t]/(t^H)$\\$Y_\nu^{m_\nu}=1$};
\node[box=rust] (splitout) at (8.4,-4.46)
  {Product over $E[t]/(t^H)$\\$Y_\nu^{m_\nu}=1$};
\draw[arrow] (in) -- node[labeltext,right] {undo the unit scaling\\lift with $z_1=t-1$} (lift);
\draw[arrow] (lift) -- node[labeltext,right] {reduce modulo $t^H$\\apply the unit scaling} (splitin);
\draw[arrow] (splitin) -- node[labeltext,above] {Fourier evaluation}
                        node[labeltext,below] {$W_S$ jet products\\and interpolation} (splitout);
\draw[arrow] (splitout) -- node[labeltext,left] {undo the unit scaling\\read $R_a(t)$ in $E[t]$} (recover);
\draw[arrow] (recover) -- node[labeltext,left] {substitute $t=1+z_1$\\truncate and scale} (out);
\end{tikzpicture}
\caption{Multiplication in a binomial block using one parameter for all
axes. The upper row uses length-$r$ jets at $t=1$ and the lower row
uses length-$H$ jets at $t=0$. Each axis wraps at most once in a
monomial product, so $H=2r+d$ captures every coefficient of the
ordinary polynomial in the middle row. This degree bound justifies
reading the polynomial before evaluating it at the other endpoint.}
\label{poly:fig:deformation}
\end{figure}


To compute these residues, pad the lifted coefficients to length $H$,
apply the isomorphism at zero, and evaluate at the $m_\nu$th roots of
unity on every axis. This identifies the split algebra with a product
of $W_S$ copies of $J_H$ because each Fourier evaluation is invertible
over $E$. Multiply pointwise, invert the tensor evaluation, and undo
the isomorphism at zero. Read the resulting length-$H$ coefficient
arrays as the actual product polynomials in $t$. Substitution of $1+z_1$
followed by truncation modulo $z_1^r$ specializes the product at the first
endpoint. The isomorphism there returns the required product in
$\mathcal T_S$. For $S=\varnothing$ the tensor evaluations are identities.
\end{proof}

For the three tensor transforms in a block $S$,
use \eqref{poly:eq:descent-sync} times $O(\log B)$, with volume
$V=W_S$, jet length $H\le3r$. The return transforms use
the same estimate as the forward transforms. The pointwise
products and the scales and shifts in the deformation lemma
also fit by Lemma~\ref{poly:lem:descent-jet-cost}
(shift by $\pm1$).
Summing over blocks proves the cyclic multiplication
bound~\eqref{poly:eq:descent-cost} in characteristic zero as well.

\begin{proof}[Proof of Theorem~\ref{thm:conv}]
Packing our ordinary convolution gives, in every characteristic and for all large $n$,
\[
 M_K(n)/n=O_\beta((1+\log T)((1+T)^\beta+\mu_K(6er))),
 \qquad \log(6er)=O(T^{4/5}\log T).
\]
For large $n$ we have $6er<n$. Induct on $n$ with the uniform
bound $M_K(n)/n\le C'(1+T)^\alpha$. The direct layer term contributes
$O_\beta((1+\log T)(1+T)^\beta)=o(T^\alpha)$.
Under the induction hypothesis the recursive term contributes at most
\[
  O_\beta\!\left((1+C')(1+\log T)
                       (1+T^{4/5}\log T)^\alpha\right)
     =(1+C')o(T^\alpha).
\]
Choose a threshold, independent of $K$ and $C'\ge1$, at which these
terms fit the claimed bound. Below it direct multiplication supplies
the base cases after increasing $C'$.
\end{proof}

\begin{corollary}[Fixed convolution and Toeplitz matrices]\label{cor:toe}
With $\alpha$ as above, any length-$n$ fixed convolution,
and hence any Toeplitz matrix on $n$ coordinates, has binary linear circuits over
the same field of size $O_\alpha(n(1+\log n)^\alpha)$.
\end{corollary}
\begin{proof}
Specialize one operand in Theorem~\ref{thm:conv} and retain the
linear parts at the gates. For a Toeplitz matrix $(t_{i-j})_{0\le i,j<n}$,
multiply the input polynomial by
$\sum_{k=-(n-1)}^{n-1}t_kX^{k+n-1}$ after padding to length $2n-1$.
The coefficient at $i+n-1$ is the desired output at $i$. These fixed
linear maps will supply the large axes in Section~\ref{sec:last}.
\end{proof}
